\par\noindent\textbf{(1)} 傍轴薄透镜的前焦面到后焦面传播实现按坐标 ξ 对应的傅里叶变换 U\_F(ξ)=C∫u(x)e\textasciicircum{}\{−i2πxξ/(λf)\}dx，常数C含相位、λ和f。频率变量 f\_x=ξ/(λf)，故傅里叶面的物理尺寸直接决定允许空间频率范围。
\par\noindent\textbf{(2)} 矩形孔径 P(ξ)=1(|ξ|≤a) 将谱乘以 H(f\_x)=1(|f\_x|≤a/(λf))。第二透镜再作反变换，像面场 u\_im(x)∝u(−x)*h(x)，其中 h(x)=∫\_\{−f\_c\}\textasciicircum{}\{f\_c\}e\textasciicircum{}\{i2πf\_xx\}df\_x=2f\_c sinc(2πf\_cx)，f\_c=a/(λf)，星号指卷积。
\par\noindent\textbf{(3)} 若物面 u(x)=u₀[1+m cos(2πfx)]，频谱为0与±f三条离散峰。当|f|<f\_c，三峰均透过，像面保留调制（坐标倒像导致相位反向）；当|f|>f\_c，±f边带被截断，仅剩均匀背景。临界|f|=f\_c需规定理想孔径边界透过约定。
\par\noindent\textbf{(4)} 有限矩形截止的h(x)具有sinc旁瓣和负振幅环，场成像可呈现振铃；这不是强度直接与h卷积，必须先算复振幅再取模方。a→∞时 f\_c→∞，h以分布意义趋δ(x)，恢复理想无截止成像。
